\title{Direct Preference Optimization: Your Language Model is Secretly a Reward Model}

\begin{abstract}
Although large unsupervised language models (LMs) acquire extensive world knowledge and certain reasoning capabilities, exerting precise control over their outputs remains challenging due to the absence of supervision during their pretraining. Contemporary approaches aiming to address this limitation typically involve collecting human preference data comparing different model outputs, then refining the unsupervised LM to better match these preferences, frequently using reinforcement learning from human feedback (RLHF). Nonetheless, RLHF involves a multi-stage and sometimes unstable workflow, which requires first learning a reward model to approximate human judgments and subsequently adjusting the original LM with reinforcement learning to optimize this learned reward, all while attempting to preserve its foundational behaviors. In this work, we present a novel parameterization for the reward model in RLHF, which allows for direct computation of the associated optimal policy in closed form. This reparameterization makes it possible to resolve the conventional RLHF objective using a straightforward classification loss function. The proposed approach, termed \textit{Direct Preference Optimization} (DPO), offers stable training, strong empirical performance, and computational efficiency, removing the need to sample from the language model during adaptation and reducing reliance on extensive hyperparameter search. Our empirical results demonstrate that DPO can adapt LMs to human preferences as effectively as, or more effectively than, previous state-of-the-art methods. In particular, DPO-based fine-tuning surpasses RLHF based on PPO in steering generation sentiment, and matches or exceeds the quality of outputs in tasks such as summarization and single-turn dialogue, all while being markedly more straightforward to implement and train.
\end{abstract}

\section{Introduction}
Unsupervised language models (LMs) of significant scale, trained on extensive text corpora, have demonstrated remarkable and sometimes unexpected proficiencies~\citep{chowdhery2022palm, brown2020language, touvron2023llama,bubeck2023sparks}. Nonetheless, the datasets employed for training such models are assembled from human-created content, reflecting diverse intentions, competencies, and preferences. Not all of these human attributes are desirable for an LM to emulate: for instance, when using an AI assistant for coding, it is beneficial for the system to \textit{comprehend} typical programming errors to facilitate corrections, yet during code generation, the LM should ideally reflect the exceptional, albeit infrequent, coding practices found in its training examples. Likewise, it is advantageous for an LM to be \textit{cognizant} of widespread misconceptions (held, for example, by half the population), but we clearly do not want the system to repeat such misconceptions as truth in proportionate numbers of interactions. Thus, distinguishing and selecting a model's \emph{preferred behaviors and responses} from its extensive \textit{learned knowledge and capabilities} is pivotal to creating AI that is robust, efficient, and amenable to control \citep{ouyang2022training}. Although standard practice involves aligning LMs to human preferences via reinforcement learning (RL), we will demonstrate that the objective typically employed in RL-based alignment can be exactly attained through a straightforward binary cross-entropy loss, thus streamlining the process of preference modeling.

Broadly, state-of-the-art techniques guide language models toward safe and useful behavior by leveraging thoughtfully curated datasets of human preferences. This preference alignment takes place subsequent to the initial unsupervised pre-training phase performed on large-scale corpora. While a direct strategy for integrating human preferences would involve supervised fine-tuning with exemplary responses authored by people, current leading approaches rely on reinforcement learning from human (or AI) feedback, known as RLHF/RLAIF~\citep{christiano2017deep,bai2022constitutional}. RLHF builds a reward model trained on human judgment data and subsequently utilizes RL algorithms to adjust a language model’s outputs to maximize these reward signals, while maintaining closeness to the base model. Despite the fact that RLHF yields systems that are highly effective in conversation and code generation, the entire RLHF workflow is far more elaborate than basic supervised methods. It necessitates training several separate language models and integrating sampling from the policy within the training loop, resulting in substantial computational overhead.

This work demonstrates a method for optimizing language models to align with human preferences directly, eliminating the need for explicit reward models or reinforcement learning frameworks. We introduce \textit{Direct Preference Optimization} (DPO), an approach that effectively targets the same optimization goal as traditional RLHF methods—namely, maximizing rewards while enforcing a KL-divergence regularization—but with a design that is both conceptually simple and easy to implement. The core idea behind the DPO update is to enhance the relative log likelihood of favored responses compared to less favored ones, while applying an adaptive, example-specific weighting mechanism that mitigates the model collapse observed with simplistic probability ratio objectives. As with prior techniques, DPO leverages a formal preference model (e.g., the Bradley-Terry model; \cite{bradley1952rankanalysis}) to evaluate how closely a reward function matches empirical preference data. Nevertheless, in contrast to established procedures, which utilize the preference model for crafting a loss to train a reward model (followed by optimizing a policy against this learned reward), DPO directly reformulates the preference loss to be a direct function of the policy through a change of variables. With access to datasets comprising human preference judgments on model outputs, DPO is able to train the policy via a straightforward binary cross-entropy loss, thereby obtaining the optimal policy under an implicit reward inferred from the observed preferences.

The primary contribution of this work is Direct Preference Optimization (DPO), an accessible algorithm that forgoes reinforcement learning to train language models from preference data. Empirical results indicate that DPO matches or exceeds the performance of contemporary approaches—including PPO-based RLHF—across various preference-driven tasks such as sentiment transfer, summarization, and dialogue, when applied to language models containing up to 6B parameters.

Increasingly large self-supervised language models have demonstrated the capacity to tackle certain tasks in a zero-shot setting \citep{radford2019language} or by employing few-shot prompt methods \citep{gpt3,megatron,chowdhery2022palm}. Yet, their effectiveness on various downstream applications and alignment with user intentions is often greatly enhanced through the process of fine-tuning on corpora containing explicit instructions and human-generated responses \citep{mishra-etal-2022-cross,sanh2022multitask,chung2022scaling,thoppilan2022lamda}. This approach, termed `instruction-tuning', enables LLMs to extrapolate to novel instructions beyond those encountered in the training set, thereby improving their general utility and flexibility \citep{chung2022scaling}. Nonetheless, collecting \textit{relative} assessments from human annotators about the quality of model outputs is frequently less demanding than gathering expert-level demonstrations. As a result, more recent studies have sought to refine LLMs using human preference datasets, yielding progress in tasks such as translation \citep{kreutzer-etal-2018-reliability}, summarization \citep{stiennon2022learning,ziegler2020finetuning}, narrative generation \citep{ziegler2020finetuning}, and following complex instructions \citep{ouyang2022training,ramamurthy2023is}. In these methodologies, a reward function parameterized by a neural network is first adapted to match the human preference data under a model like Bradley-Terry \citep{bradley1952rankanalysis}, after which the language model undergoes further optimization to maximize this reward, utilizing reinforcement learning strategies such as REINFORCE \citep{williams1992reinforce}, proximal policy optimization (PPO; \cite{schulman2017proximal}), or similar techniques \citep{ramamurthy2023is}. Parallel research leverages LLMs tuned with both instruction-following capabilities and human feedback to autonomously generate additional synthetic preference data for specific traits like safety or non-harmfulness \citep{bai2022constitutional}, relying merely on lightly-supervised text-based rubrics guiding the LLM's annotations. Collectively, these advances represent a synthesis of two lines of inquiry: one focusing on the application of reinforcement learning to train language models across a spectrum of tasks~\citep{Ranzato2015SequenceLT,paulus2018a,wu2018learning}, and the other developing broad frameworks for learning from human preferences \citep{christiano2017deep,kupcsik2018learning}. While learning from relative human preferences offers conceptual and practical appeal, implementing reinforcement learning-based fine-tuning in large-scale language models is still highly nontrivial; the present work addresses this with a theoretically-grounded alternative for optimizing relative preference signals that obviates the need for RL.

Research on learning policies from preferences outside linguistic tasks has been explored within both bandit and reinforcement learning frameworks, leading to a variety of methods. In particular, the contextual dueling bandit (CDB; \cite{yue2012karmed,dudik2015contextual}) paradigm addresses contextual bandit problems by leveraging preference or ranking information over actions in place of scalar rewards. In this scenario, since absolute rewards are unavailable, theoretical investigations in CDBs introduce the concept of a \textit{von Neumann winner}: this refers to a policy whose average likelihood of prevailing over \textit{any} competing policy is no less than 50\% \citep{dudik2015contextual}. Notably, CDB frameworks acquire preference annotations in an online manner, whereas learning from human preferences typically utilizes a static, offline collection of action pairs labeled with preferences \citep{yan2022human}. Analogously, \textit{preference-based RL} (PbRL) studies settings in which feedback is delivered as binary preferences provided by an \textit{unknown} `scoring' or reward function, as opposed to explicit rewards \citep{BusaFekete2014,ruiz2023dueling}. A number of PbRL techniques have been introduced—including methods capable of making use of off-policy preference data—but most of them follow a two-stage process: they begin by estimating the underlying (latent) scoring or reward model, and subsequently optimize policies with respect to this estimated reward \citep{jain2013learning,BusaFekete2014,christiano2017deep,sadigh2017active,kupcsik2018learning}. In contrast, we propose a unified approach in which policy learning is carried out in a single stage by directly optimizing policies to fulfill the specified preferences.

\section{Preliminaries}\label{section:prelims}

We now revisit the RLHF pipeline as established in \citeauthor{ziegler2020finetuning} and further elaborated by works such as \citep{stiennon2022learning, bai2022training, ouyang2022training}. This pipeline generally consists of three components: (1) supervised fine-tuning (SFT); (2) acquisition of preferences and subsequent reward modelling; and (3) reinforcement learning optimization.

\textbf{SFT}: The standard RLHF process commences by adapting a pre-trained language model to downstream applications (e.g., dialogue systems, text summarization) using supervised learning on carefully curated datasets, resulting in the model $\pi^\text{SFT}$.

\textbf{Reward Modelling Phase}: In the subsequent step, prompts $x$ are used to elicit response pairs $(y_1, y_2)\sim \pi^\text{SFT}(y \mid x)$ from the SFT model. Human evaluators are then presented with these pairs and asked to indicate their preferred response, which is denoted as $y_w\succ y_l \mid x$; here, $y_w$ and $y_l$ refer to the favored and less-favored completions within $(y_1, y_2)$. It is assumed that these preference judgments are produced by an implicit, unobserved reward function $r^*(y, x)$. Among several frameworks for modeling preferences, the Bradley-Terry (BT) model \cite{bradley1952rankanalysis} is commonly adopted, although more general models—such as the Plackett-Luce family \citep{plackett1975analysis, luce2012individual}—are applicable when multiple ranked options are present. Under the BT formulation, the distribution over human preferences $p^*$ takes the following form:

\begin{equation}\label{eq:bradley-terry}
    p^*(y_1\succ y_2 \mid x)=\frac{\exp\left(r^*(x, y_1)\right)}{\exp\left(r^*(x, y_1)\right) + \exp\left(r^*(x, y_2)\right)}.
\end{equation}

Given a fixed dataset of pairwise comparisons $\mathcal{D}=\bigl\{x^{(i)}, y_w^{(i)}, y_l^{(i)}\bigr\}_{i=1}^N$ drawn from the distribution $p^*$, one can introduce a parameterized reward function $r_{\phi}(x, y)$ and learn its parameters by maximizing the likelihood. When framed as a binary classification task, the corresponding objective is a negative log-likelihood loss:

\begin{equation}\label{eq:reward_model}
    \mathcal{L}_R(r_{\phi}, \mathcal{D}) = -\mathbb{E}_{(x, y_w, y_l)\sim \mathcal{D}}\bigl[\log \sigma(r_{\phi}(x, y_w)- r_{\phi}(x, y_l))\bigr]
\end{equation}

where $\sigma$ denotes the logistic sigmoid function. In large language models, $r_{\phi}(x, y)$ typically reuses the pretrained SFT model $\pi^\text{SFT}(y \mid x)$, with a linear layer added atop the last transformer output, producing a single real-valued reward prediction \cite{ziegler2020finetuning}. To limit the variance of the resulting reward estimates, prior research commonly applies normalization such that $\mathbb{E}_{x,y\sim \mathcal{D}}\left[r_\phi(x, y)\right] = 0$ for every $x$.

\textbf{RL Fine-Tuning Phase}: In the subsequent RL fine-tuning phase, the trained reward model is used to supply evaluative signals to the language model. Following established practice~\citep{jaques2017sequence, jaques2020human}, the optimization goal takes the form

\begin{equation}\label{eq:RL}
\max_{\pi_{\theta}}  \mathbb{E}_{x\sim \mathcal{D}, y\sim \pi_{\theta}(y \mid x)}\bigl[r_{\phi}(x, y)\bigr] - \beta\mathbb{D}_{\textrm{KL}}\bigl[\pi_{\theta}(y\mid x)\mid \mid \pi_\text{ref}(y\mid x)\bigr],
\end{equation}

where $\beta$ modulates the strength of regularization relative to the reference policy $\pi_\text{ref}$—often, this reference is the original SFT model $\pi^\text{SFT}$. In implementation, the new policy $\pi_\theta$ is generally initialized with weights from $\pi^\text{SFT}$. The KL-divergence penalty is crucial: it restricts the model’s outputs from departing too much from the region where the reward model yields trustworthy estimates, supports response diversity, and helps mitigate the risk of the policy collapsing to only a few answers with high rewards. Due to the discrete nature of sequence generation, the overall objective lacks differentiability and is hence commonly handled with reinforcement learning algorithms. The prevailing method \citep{ziegler2020finetuning, stiennon2022learning, bai2022training, ouyang2022training} is to define the reward for policy optimization as ${r(x, y) = r_{\phi}(x, y) -\beta (\log \pi_{\theta}(y\mid x) - \log \pi_\text{ref}(y\mid x))}$ and optimize it using PPO \cite{schulman2017proximal}.

\section{Direct Preference Optimization}\label{sec:DPO}

Prompted by the complexity of deploying reinforcement learning at scale for language model fine-tuning, we seek a streamlined technique for policy optimization based directly on preferences. In contrast with conventional RLHF strategies—which require learning a reward function before RL-based policy tuning—our method relies on a specific reward model formulation, from which one can derive its optimal policy analytically, eliminating the need for a reinforcement learning stage. As detailed below, our principal contribution is to employ a closed-form link between the reward function and the corresponding optimal policy. This enables reparameterizing a loss over rewards as a loss over policies, thereby removing the need to fit a distinct reward model, while still aligning with human preference models such as the Bradley-Terry framework. Effectively, the policy

\textbf{Deriving the DPO objective.} We begin by considering the reinforcement learning objective employed in previous literature, as stated in Eq.~\ref{eq:RL}, for an arbitrary reward function $r$. Building on prior studies~\citep{peters2007reinforcement, peng2019advantage, korbak2022reinforcement, go2023aligning}, it can be readily demonstrated that the optimal solution to the KL-regularized reward maximization problem defined in Eq.~\ref{eq:RL} is characterized by:

\begin{equation}\label{eq:op_policy}
    \pi_r(y\mid x) = \frac{1}{Z(x)}\pi_\text{ref}(y\mid x)\exp\left(\frac{1}{\beta}r(x, y)\right),
\end{equation}

where the partition function $Z(x)$ is given by $Z(x) =\sum_{y}\pi_\text{ref}(y\mid x)\exp\left(\frac{1}{\beta}r(x, y)\right)$. A detailed derivation is provided in Appendix \ref{app:derivation1}. Direct computation of $Z(x)$ remains computationally demanding, even when employing the MLE estimate $r_{\phi}$ as a surrogate for the ground-truth reward $r^*$~\citep{korbak2022reinforcement, go2023aligning}, presenting practical challenges. Nevertheless, Eq.~\ref{eq:op_policy} can be rearranged to re-express the reward in terms of its associated optimal policy $\pi_r$, the reference policy $\pi_\text{ref}$, and the partition function $Z(\cdot)$. By taking the logarithm of both sides of Eq.~\ref{eq:op_policy} and simplifying, we have:

\begin{equation}\label{eq:main_eq}
    r(x,y) =\beta \log \frac{\pi_r(y\mid x)}{\pi_\text{ref}(y\mid x)} + \beta \log Z(x).
\end{equation}

This formulation allows for substitution of the true reward function $r^*$ and its corresponding optimal policy $\pi^*$. Importantly, the Bradley-Terry model is dependent only upon the difference in rewards for a pair of outputs, i.e., ${p^*(y_1 \succ y_2 \mid x) = \sigma(r^*(x, y_1) - r^*(x, y_2))}$. If we now substitute the reparameterization from Eq.~\ref{eq:main_eq} for $r^*(x,y)$ into the Bradley-Terry preference model of Eq.~\ref{eq:bradley-terry}, the partition function terms eliminate, leading to a human preference probability that depends solely on the optimal policy $\pi^*$ and the reference policy $\pi_\text{ref}$. Hence, under the Bradley-Terry model, the optimal RLHF policy $\pi^*$ satisfies the following preference expression:

\begin{equation}\label{eq:objective}
    p^*(y_1\succ y_2 \mid x)=\frac{1}{1 + \exp\left(\beta \log \frac{\pi^*(y_2\mid x)}{\pi_\text{ref}(y_2\mid x)} - \beta \log \frac{\pi^*(y_1\mid x)}{\pi_\text{ref}(y_1\mid x)}\right)}
\end{equation}

The full details of this derivation are presented in Appendix~\ref{app:derivation2}. While Eq.~\ref{eq:objective} specifically utilizes the Bradley-Terry model, comparable derivations are available for more general cases such as the Plackett-Luce family~\citep{plackett1975analysis, luce2012individual}, as documented in Appendix~\ref{app:plackett_luce_models}.

Having formulated the likelihood of the human preference data in terms of the optimal policy, we can now define a maximum likelihood objective with respect to a parameterized policy $\pi_\theta$. Similar in spirit to the reward-based formulation (see Eq.~\ref{eq:reward_model}), our policy-level objective takes the form:

\begin{equation}\label{eq:optimum_model}
    \mathcal{L}_\text{DPO}(\pi_{\theta}; \pi_\text{ref}) = -\mathbb{E}_{(x, y_w, y_l)\sim \mathcal{D}}\left[\log \sigma \left(\beta \log \frac{\pi_{\theta}(y_w\mid x)}{\pi_\

Through this approach, we learn an implicit reward via an alternative parameterization in which the optimal policy coincides with $\pi_\theta$. Additionally, our method corresponds to fitting a reparameterized Bradley-Terry model, thereby inheriting certain theoretical guarantees, including consistency under appropriate conditions on the distribution of preference data \cite{bong2022generalized}. Section~\ref{sec:theory} offers further discussion on the theoretical aspects of DPO and its relationship to existing literature.

\textbf{Mechanistic Interpretation of the DPO Update.} To gain intuition about the behavior of DPO, it is instructive to inspect the gradient of its objective function, $\mathcal{L}_\text{DPO}$. The derivative with respect to the model parameters $\theta$ can be expressed as follows:

\begin{multline*}\label{eq:gradient}
    \nabla_\theta \mathcal{L}_\text{DPO}(\pi_\theta;\pi_\text{ref}) = \\ -\beta\mathbb{E}_{(x, y_w, y_l) \sim \mathcal{D}} \bigg[\underbrace{\sigma(\hat{r}_\theta(x, y_l) - \hat{r}_\theta (x, y_w))}_\text{greater emphasis when reward ranking is incorrect}\bigg[\underbrace{\nabla_\theta\log \pi(y_w \mid x)}_\text{promote $y_w$'s likelihood} - \underbrace{\nabla_\theta\log\pi(y_l \mid x)}_\text{diminish $y_l$'s likelihood}\bigg]\bigg],
\end{multline*}

where $\hat{r}_\theta(x, y) = \beta \log \frac{\pi_\theta(y \mid x)}{\pi_\text{ref}(y \mid x)}$ denotes the reward inferred from the language model $\pi_\theta$ relative to the reference $\pi_\text{ref}$ (see Section~\ref{sec:theory} for details). Conceptually, this gradient drives up the probability of favored outputs $y_w$ while suppressing that of less favored ones $y_l$. Notably, the influence of each training instance is modulated by the degree to which the implicit reward function $\hat{r}_\theta$ erroneously ranks the dispreferred completion higher, scaled by $\beta$; this reflects the extent to which the reward function misorders the pair, modulated by the KL divergence penalty. Empirical findings indicate this weighting is crucial—removing it, as in a naïve implementation, can lead the language model to pathological behavior (see Appendix Table~\ref{tab:unlikelihood_generations}).

\textbf{DPO overview.} The typical DPO process comprises: 1) Drawing completion samples $y_1, y_2 \sim \pi_\text{ref}(\cdot \mid x)$ for each prompt $x$, then annotating these pairs with human preference data to create the preference dataset $\mathcal{D} = \{x^{(i)}, y_w^{(i)}, y_l^{(i)}\}_{i=1}^N$; 2) Training the language model $\pi_\theta$ by minimizing $\mathcal{L}_\text{DPO}$, utilizing the fixed $\pi_\text{ref}$, dataset $\mathcal{D}$, and chosen $\beta$. 
In applied scenarios, it is generally preferable to leverage existing publicly released preference datasets instead of generating new model outputs and gathering additional human preference labels. Because such datasets are typically derived from sampling with $\pi^\text{SFT}$, we choose $\pi_\text{ref} = \pi^\text{SFT}$ whenever this is possible. If no $\pi^\text{SFT}$ is provided, we instead set $\pi_\text{ref}$ by fitting it via maximum likelihood estimation to the preferred responses ${(x, y_w)}$, specifically: ${\pi_\text{ref} = \argmax_{\pi}\mathbb{E}_{x, y_w \sim \mathcal{D}}\left[\log \pi(y_w \mid x)\right]}$. Employing this initialization procedure reduces the mismatch between the unavailable ideal reference distribution and the practical $\pi_\text{ref}$ utilized by DPO. Appendix~\ref{app:implementation} provides additional information regarding implementation choices and hyperparameters.

\section{Theoretical Analysis of DPO} This section offers a deeper theoretical understanding of DPO, clarifying its mechanisms, establishing formal justification, and articulating the benefits of DPO over actor-critic approaches such as those used in RLHF (for example, PPO~\cite{schulman2017proximal}).

\label{sec:theory}

\subsection{Your Language Model Is Secretly a Reward Model} The DPO framework enables policy learning through a single maximum likelihood formulation, thereby sidestepping the need for both explicit reward function estimation and reinforcement learning steps. The loss in Eq. \ref{eq:main_eq} aligns with the Bradley-Terry model, utilizing the reward parameterization $r^*(x, y) = \beta \log\frac{\pi^*_\theta(y \mid x)}{\pi_\text{ref}(y \mid x)}$, and optimizing the parameterized model $\pi_{\theta}$ is therefore analogous to the reward model training described in Eq. \ref{eq:reward_model}, modulo a change of variables. In what follows, we formalize the underpinning theory of this parameterization, demonstrate that it places no undue restrictions on the family of learnable reward models, and prove that it admits precise recovery of the optimal policy. We begin by introducing an equivalence relation over reward functions.

\begin{definition}
Two reward functions $r(x, y)$ and $r'(x, y)$ are called equivalent if and only if there exists a function $f$ such that ${r(x, y)-r'(x, y) = f(x)}$.
\end{definition}

It is straightforward to verify that this defines an equivalence relation, which organizes the set of all reward functions into equivalence classes. The following lemmas formalize two important consequences of this equivalence:

\begin{lemma}\label{lemma:same_prefrence}
Within the Plackett-Luce framework, and specifically under the Bradley-Terry model, any pair of reward functions belonging to the same equivalence class will yield identical preference distributions.
\end{lemma}

\begin{lemma}\label{lemma:same_policy}
    Any two reward functions from the same equivalence class lead to the same optimal policy for the constrained RL problem.
\end{lemma}

The proofs are elementary and can be found in Appendix \ref{app:lemma1}. The first lemma highlights a well-known identifiability issue associated with the Plackett-Luce family of models \cite{plackett1975analysis}, resulting in under-specification. Consequently, additional constraints must often be imposed to ensure identifiability when estimating parameters via maximum likelihood (see Eq.~\ref{eq:reward_model}) \cite{bong2022generalized}. According to the second lemma, all reward functions in the same class produce the same optimal policy, which justifies focusing on recovering any single representative from the correct equivalence class. The main theorem, whose proof is deferred to Appendix~\ref{app:thm1}, establishes the following:

\begin{theorem}\label{thm:main}
    Provided mild conditions hold, every reward class that is consistent with the Plackett-Luce (and, specifically, the Bradley-Terry) model admits a representation of the form ${r(x, y) = \beta \log \frac{\pi(y\mid x)}{\pi_\text{ref}(y\mid x)}}$ for a suitable model $\pi(y\mid x)$ and a given reference distribution $\pi_\text{ref}(y \mid x)$.
\end{theorem}

\begin{sproof}
    Let $r(x, y)$ be an arbitrary reward function with corresponding optimal model $\pi_r(y \mid x)$ as defined in Eq.~\ref{eq:op_policy}. We claim that a member of the equivalence class of $r$ can always be written via the proposed reparameterization. Consider the following projection:
\begin{equation}
    f(r; \pi_\text{ref}, \beta)(x, y) = r(x, y) - \beta\log\sum_{y}\pi_\text{ref}(y\mid x)\exp\left(\frac{1}{\beta}r(x, y)\right)
\end{equation}
This operator $f$ acts to normalize $r(x, y)$ using the log-partition function induced by $\pi_r$. The added normalization term depends solely on $x$, ensuring $f(r; \pi_\text{ref}, \beta)(x, y)$ remains in the same equivalence class as $r(x, y)$. Substituting the right-hand side of Eq.~\ref{eq:main_eq} for $r$, valid for any reward function, gives $f(r; \pi_\text{ref}, \beta)(x, y) = \beta \log \frac{\pi_r(y\mid x)}{\pi_\text{ref}(y\mid x)}$. Thus, this projection yields a function of the desired form within the equivalence class of $r$, establishing that our reparameterization is general enough for the reward modeling framework considered.
\end{sproof}

An alternative interpretation of Theorem~\ref{thm:main} is that it precisely characterizes the particular reward function, from each equivalence class, that is selected by the DPO reparameterization—namely, the reward function that fulfills the following condition:

\begin{equation}\label{eq:lag_p}
     \sum_{y}\underbrace{\pi_\text{ref}(y\mid x)\exp\left(\frac{1}{\beta}r(x, y)\right)}_{=\pi(y\mid x)\text{, via Thm.~\ref{thm:main} reparam.}} = 1,
\end{equation}

which ensures that $\pi(y\mid x)$ forms a legitimate probability distribution (i.e., all probabilities are nonnegative and sum to unity). Examining Eq.~\ref{eq:op_policy} reveals that Eq.~\ref{eq:lag_p} constitutes the partition function corresponding to the optimal policy derived from the reward function $r(x, y)$. The central insight driving the DPO algorithm lies in its ability to introduce specific constraints within the otherwise underdetermined Plackett-Luce family (notably including the Bradley-Terry model) of preference-based frameworks. This approach maintains the expressive capacity of the reward models, yet guarantees that the optimal policy as given by Eq.~\ref{eq:op_policy} remains analytically manageable for arbitrary prompts $x$.

\subsection{Instability of Actor-Critic Algorithms} Our framework can additionally elucidate sources of instability common to traditional actor-critic approaches in RLHF contexts, such as PPO. In alignment with the RLHF methodology, we focus here on the RL fine-tuning stage discussed in Section \ref{section:prelims}. By making parallels to the control as inference paradigm \cite{levine2018reinforcement} in the setting of constrained reinforcement learning specified by \ref{eq:RL}, we suppose the policy model $\pi_{\theta}(y\mid x)$ is parameterized and optimize the objective $\mathbb{D}_{\text{KL}}[\pi_{\theta}(y|x) \mid \mid \pi^*(y\mid x)]$, with $\pi^*$ denoting the optimal policy associated with the reward function $r_{\phi}(y, x)$ as described in Eq.~\ref{eq:optimum_model}. Straightforward manipulations yield the following objective function:

\begin{equation}\label{eq:AC}
    \max_{\pi_{\theta}}\mathbb{E}_{\pi_{\theta}(y\mid x)}\bigg[\underbrace{r_{\phi}(x, y) -\beta\log\sum_{y}\pi_\text{ref}(y\mid x)\exp\left(\frac{1}{\beta}r_{\phi}(x, y)\right)}_{f(r_{\phi}, \pi_\text{ref}, \beta)} - \underbrace{\beta\log\frac{\pi_{\theta}(y\mid x)}{\pi_\text{ref}(y\mid x)}}_{\text{KL}}\bigg]
\end{equation}

The objective here aligns with that optimized in previous studies 
\citep{ziegler2020finetuning, stiennon2022learning, bai2022training, ouyang2022training}, specifically utilizing the DPO-equivalent reward for the reward class $r_{\phi}$. In this context, the normalization component within $f(r_{\phi}, \pi_\text{ref}, \beta)$ can be interpreted as representing the soft value function of the reference policy $\pi_\text{ref}$. Although this normalization does not influence the ultimate optimality of solutions, omitting it may cause the policy gradient estimator to exhibit substantial variance, potentially compromising training stability. Incorporating this normalization via a learned value function is an option, but presents its own optimization challenges. Another method, explored in prior work, involves reward normalization based on a human completion baseline, effectively a single-sample Monte Carlo estimate for normalization. By contrast, the DPO reparameterization results in a reward formulation that dispenses with the need for any such baseline. 

\section{Experiments}
This section presents empirical investigations of DPO's effectiveness in learning policies from preference data. We begin by probing, within a controlled text generation framework, how well DPO balances the objectives of reward maximization and minimal KL-divergence to the reference policy in comparison to prevalent preference-based algorithms such as PPO. Subsequently, DPO is assessed on larger-scale models and more challenging RLHF tasks, notably in summarization and conversational domains. Our findings reveal that with negligible hyperparameter adjustment, DPO generally matches or surpasses established methods—including RLHF with PPO and the best-of-$N$ selection from multiple sampled trajectories under a trained reward model. We first outline the experimental methodology, with further specifics provided in Appendix~\ref{app:exp_details}.

\textbf{Tasks.} We investigate three distinct open-ended text generation tasks in our study. Throughout all experiments, the various algorithms are trained to learn a policy using a dataset of preferences, denoted by $\mathcal{D}=\bigl\{x^{(i)}, y_w^{(i)}, y_l^{(i)}\bigr\}_{i=1}^N$. In the case of \textbf{controlled sentiment generation}, the input $x$ comprises the initial segment of a movie review sourced from the IMDb dataset \cite{maas-EtAl:2011:ACL-HLT2011}, and the aim is for the policy to produce a continuation $y$ reflecting positive sentiment. To facilitate controlled comparison, preference pairs for this task are automatically constructed over model generations utilizing a pre-trained sentiment classifier, ensuring that $p(\text{positive}\mid x,y_w)>p(\text{positive}\mid x,y_l)$. For SFT, we perform continued training of GPT-2-large until convergence, using reviews from the training portion of the IMDb dataset (see App~\ref{app:sentiment_details} for additional information). For the \textbf{summarization} task, $x$ is a Reddit forum post, with the policy responsible for producing a summary $y$ that encapsulates the main points. Consistent with prior literature, we rely on the Reddit TL;DR summarization dataset \citep{volske-etal-2017-tl} and incorporate human preference data assembled by \citeauthor{stiennon2022learning}. The SFT model here is trained on summaries written by humans\footnote{\url{https://huggingface.co/CarperAI/openai_summarize_tldr_sft}} using the TRLX \citep{leandro_von_werra_2023_7790115} toolkit for RLHF. Notably, the preference dataset derives from samples produced by a different, though comparably trained, SFT model as gathered by \citeauthor{stiennon2022learning}. Lastly, in the \textbf{single-turn dialogue} scenario, $x$ consists of a user's query—ranging from complex questions in astrophysics to personal advice requests. Here, the policy generates an informative and engaging reply $y$; we make use of the Anthropic Helpful and Harmless dialogue corpus \citep{bai2022training}, encompassing 170,000 exchanges between humans and automated agents. Each dialogue culminates in two model-generated candidate responses, one of which is labeled as preferred by a human annotator. Since a ready-made SFT model is unavailable in this context, we derive the SFT model by fine-tuning a publicly available language model solely on the preferred responses.

\textbf{Evaluation.} We adopt two primary evaluation strategies in our experiments. For the controlled sentiment generation task, where the true reward function (as given by a sentiment classifier) is available, we assess each algorithm by plotting its Pareto frontier of achieved reward versus KL-divergence from the reference policy. This allows for a direct comparison with respect to the constrained reward maximization objective. By contrast, in realistic scenarios where the true reward function is unavailable, we employ a \textit{win rate} metric to measure the performance of each algorithm relative to a baseline policy. For this, we rely on GPT-4 to serve as a surrogate for human judgment, evaluating either the quality of summaries (summarization) or the helpfulness of responses (single-turn dialogue). The baseline for summarization comprises the reference summaries from the test set, while for dialogue, it is the human-preferred response in the test data. Although prior research indicates that language models can outperform traditional metrics in automated evaluation \citep{Chen2023ExploringTU}, we further support the use of GPT-4 in our evaluation framework through a human study presented in Sec.~\ref{sec:human-judgments}. The outcomes indicate a strong correlation between GPT-4 assessments and human judgments, with the consistency between humans and GPT-4 matching or surpassing inter-annotator agreement levels.

\textbf{Methods.} Alongside DPO, we benchmark a variety of established methods for aligning language models with human preferences. For the summarization task, we use zero-shot prompting with \textbf{GPT-J} \citep{gpt-j}, while for dialogue we implement 2-shot prompting with \textbf{Pythia-2.8B} \citep{biderman2023pythia}. We also include evaluations of the \textbf{SFT} model and the \textbf{Preferred-FT} variant, which is obtained through supervised fine-tuning on the selected completion $y_w$, sourced either from the SFT model (in sentiment and summarization) or from a general language model (in the case of dialogue). Another comparison method, \textbf{Unlikelihood}~\citep{welleck2019neural}, modifies training to increase the probability of $y_w$ while explicitly decreasing the probability of $y_l$, modulated by an optional parameter $\alpha\in[0,1]$ weighting the unlikelihood component. Furthermore, we assess \textbf{PPO} \citep{schulman2017proximal} trained with a reward model learned from preference data, as well as \textbf{PPO-GT}, which acts as an oracle trained using the actual reward function in the controlled sentiment context. For sentiment tasks, PPO-GT is evaluated in two forms: a standard implementation \cite{leandro_von_werra_2023_7790115} and an enhanced version with reward normalization and hyperparameter optimization, enhancements we also adopt for PPO with learned rewards. Finally, our baselines include the \textbf{Best of $N$} strategy: here, $N$ responses are generated from either the SFT or Preferred-FT model (depending on the task), and the response with the highest score, according to a preference-trained reward model, is selected. While this approach can yield strong performance, it is not computationally scalable for even moderate $N$ since it requires generating $N$ completions for each query during inference.

\subsection{How well can DPO optimize the RLHF objective?}

In standard RLHF methods, the KL-constrained reward maximization objective is designed to balance maximizing the reward signal against keeping the learned policy close to the reference policy, as measured by KL divergence. Thus, an appropriate comparison of algorithms must jointly consider both the attained reward and the corresponding KL value; higher reward accompanied by a much larger KL may not constitute a true improvement. Figure~\ref{fig:frontier-tldr-main} displays the frontier of achievable reward versus KL values for a selection of algorithms in the sentiment analysis task. For each algorithm, we perform several training runs with different policy conservativeness settings (target KL $\in\{3,6,9,12\}$ for PPO, $\beta \in \{0.05,0.1,1,5\}$, $\alpha\in\{0.05,0.1,0.5,1\}$ for unlikelihood, and varying random seeds for preferred-FT), amounting to a total of 22 experimental runs. After every 100 training steps up to convergence, we assess the learned policies by evaluating them on a fixed test prompt set, calculating both the mean reward according to the ground-truth reward function and the mean sequence-level KL\footnote{Specifically, we compute the sum of KL divergences over all timesteps of the sequence.} with respect to the reference policy $\text{KL}\left(\pi\mid \mid \pi_\text{ref}\right)$. Our results indicate that DPO yields the most favorable frontier, securing high rewards with relatively low KL. This is noteworthy for several reasons. First, despite both DPO and PPO aiming to optimize the same underlying objective, DPO consistently exhibits greater efficiency, achieving a strictly better reward/KL balance than PPO. Second, DPO outperforms PPO even in scenarios where PPO is given access to true rewards (PPO-GT), surpassing its reward-KL frontier.

\subsection{Can DPO scale to real preference datasets?}
\label{sec:dpo-real-datasets}
We proceed to assess the fine-tuning capabilities of DPO on tasks involving summarization and single-turn dialogue. It is well-established that for summarization, standard automated metrics like ROUGE often fail to reliably align with human judgments~\citep{stiennon2022learning}. Prior research indicates that preference-based fine-tuning of language models with PPO can yield summaries that better reflect user preferences. To evaluate various approaches, we generate model outputs using the test partition of the TL;DR summarization dataset, measuring the average win rate of these completions relative to the gold-standard references in the test data. Samples are drawn from all methods using a range of temperature values between 0.0 and 1.0, with corresponding win rates visualized in Figure~\ref{fig:frontier-tldr-main} (right). DPO, PPO, and Preferred-FT models are all trained by further fine-tuning the same GPT-J SFT model\footnote{\url{https://huggingface.co/CarperAI/openai_summarize_tldr_sft}}. Our results show that at a temperature of 0.0, DPO achieves a win rate of about 61\%, outperforming PPO, which achieves roughly 57\% at its optimal (0.0) temperature setting. Additionally, DPO's peak win rate surpasses that of the best-of-$N$ baseline approach. It is worth mentioning that no significant hyperparameter tuning was applied to DPO's $\beta$, so these findings might not reflect the upper bounds of DPO's capabilities. Further, DPO demonstrates greater resilience to changes in sampling temperature, whereas PPO's performance deteriorates substantially at higher temperatures, approaching that of the original GPT-J model. Preferred-FT exhibits negligible improvements over the SFT baseline. Lastly, in our head-to-head human evaluation (see Section~\ref{sec:human-judgments}), DPO outputs sampled at temperature 0.25 are preferred 58\% of the time over PPO outputs generated at temperature 0.

For single-turn dialogue tasks, we assess various approaches on the subset of the Anthropic HH dataset \citep{bai2022training} test split that contains only one round of interaction between a human and an assistant. To determine win rates for each method, GPT-4 evaluations use the preferred test set completions as the ground truth reference. In the absence of a standardized SFT model for this evaluation, our process begins with a pre-trained Pythia-2.8B model; we first apply Preferred-FT to fine-tune a reference model on selected completions so that its outputs remain consistent with the training distribution, followed by training with DPO. Additionally, we benchmark against the top completion chosen from 128 Preferred-FT generations (as the Best of $N$ method shows diminishing returns beyond 128 completions for this problem—see Appendix Figure~\ref{fig:best-of-n}) and a 2-shot prompt version of the base Pythia-2.8B model, noting that DPO matches or surpasses the best-performing temperatures for each technique. We further assess an RLHF system using PPO, specifically trained on the Anthropic HH dataset \footnote{\url{https://huggingface.co/reciprocate/ppo_hh_pythia-6B}}, sourced from a reputable implementation \footnote{\url{https://github.com/CarperAI/trlx/tree/main/examples/hh}}, but are unable to identify prompt settings or temperature values that yield results superior to those of the original Pythia-2.8B base model. In light of findings from TL;DR and given that both PPO and Best of 128 optimize identical reward criteria, we treat Best of 128 as a practical stand-in for PPO-level performance. In sum, DPO is unique among computationally lightweight techniques in improving upon the preferred completions in the Anthropic HH set and offers performance on par with or exceeding the more resource-intensive Best of 128 baseline. Finally, as depicted in Figure~\ref{fig:dialogue-main}, DPO attains optimal performance after relatively few training steps.

\subsection{Generalization to a new input distribution}

To better assess how PPO and DPO perform under distributional changes, we analyze the policies trained in our Reddit TL;DR summarization task by deploying them on a distinct data distribution: the test set of the CNN/DailyMail news articles \citep{nallapati-etal-2016-abstractive}. The evaluations employ the optimal sampling temperatures derived from TL;DR (0 and 0.25). Table~\ref{tab:ood} presents these outcomes. For comparison, we determine the win rate of each policy versus reference summaries in the test data using GPT-4 and the identical GPT-4 (C) prompt from the Reddit TL;DR evaluation, substituting “forum post” with “news article.” Notably, on this shifted distribution, DPO consistently exceeds the performance of PPO by a notable margin. These results suggest that DPO exhibits a comparable capacity for distributional generalization as PPO, even though DPO does not benefit from the unlabeled Reddit TL;DR prompts leveraged by PPO.

\subsection{Validating GPT-4 judgments with human judgments}
\label{sec:human-judgments}

To ascertain the fidelity of GPT-4 as an evaluator, we conduct a human rating study using the outputs from the TL;DR summarization experiment, assessing with two GPT-4 prompts. The first prompt, \textbf{GPT-4 (S)} (simple), asks GPT-4 to judge which summary best captures the critical information from a post. The second, \textbf{GPT-4 (C)} (concise), extends this by requesting that the preferred summary also be more concise—this prompt is of interest since GPT-4 appears biased toward lengthier, more repetitive summaries when given the \textbf{GPT-4 (S)} instruction, diverging from human raters. Full prompt wording can be found in Appendix~\ref{app:prompts}. Our study compares three summarization systems: the strongest (DPO at temp. 0.25), the weakest (PPO at temp. 1.0), and an intermediate performer (SFT at temp. 0.25), with each assessed in head-to-head comparisons against the top-performing greedily-decoded PPO system. Across both prompts, we observe that GPT-4’s judgments are as consistent with human ratings as human raters are with each other, indicating that GPT-4 serves as an adequate stand-in for human evaluators (due to constraints on rater availability, repeated ratings were only collected for the DPO and PPO-1 comparisons). On the whole, the \textbf{GPT-4 (C)} prompt yields win rates more closely aligned with human preferences, which motivates our use of this prompt in Section~\ref{sec:dpo-real-datasets}. Further methodology and information about rater instructions and the participant pool are available in Appendix~\ref{app:human-study}.

\section{Discussion}
Preference-based learning constitutes a robust and scalable method for developing language models that are both competent and aligned. In this work, we presented DPO, a streamlined training procedure that leverages preference data to train language models, eliminating the need for reinforcement learning. Instead of reformulating the preference learning task to conform to traditional RL paradigms to make use of standard RL algorithms, DPO establishes a correspondence between language model policies and reward functions. This allows for direct training of language models to adhere to human preferences using a straightforward cross-entropy loss, maintaining generality and bypassing reinforcement learning. Remarkably, with minimal hyperparameter adjustment, DPO achieves performance on par with, or surpassing, existing RLHF methods, such as those using PPO, thereby significantly lowering the complexity and resource demands of preference-based model training.

\textbf{Limitations \& Future Work.} The outcomes of this research suggest multiple avenues for further investigation. One important question concerns how well DPO-trained policies extrapolate beyond their training distributions compared to approaches that utilize explicit reward functions; preliminary evidence indicates comparable generalization to models trained via PPO, but thorough analysis remains necessary. Additionally, it is an open question whether leveraging self-labeling mechanisms within DPO could enable effective utilization of prompts lacking human annotation. Another line of inquiry involves understanding the dynamics of reward over-optimization within the DPO framework, particularly whether the modest performance decline observed in Figure~\ref{fig:dialogue-main}-right exemplifies this phenomenon. While our assessment encompasses models with up to 6B parameters, extending DPO to significantly larger, state-of-the-art architectures represents a promising area for future work. On the evaluation front, our findings indicate that GPT-4’s computed win rates are sensitive to prompt formulation; future research might explore strategies for obtaining higher-fidelity judgments from automated evaluators. Beyond the domain of language model training, DPO offers considerable potential for broader application, such as in the optimization of generative models across various data modalities.